
   \begin{longtable}[c]{|p{0.7cm}|p{1.9cm}|p{1.9cm}|p{1.9cm}|p{1.9cm}|p{1.9cm}|p{1cm}|p{1cm}|p{1.2cm}| }
   \caption{IO MUX Functions List}
    \label{table:iomux_pads} \\
        \hline
        \rowcolor{lightgray}
\textbf{GPIO} & \textbf{Pin Name} & \textbf{Function 0} & \textbf{Function 1} & \textbf{Function 2} & \textbf{Function 3} & \textbf{DRV} & \textbf{Reset} & \textbf{Notes} \\
 \hline
\endfirsthead
        \hline
        \rowcolor{lightgray}
        \textbf{GPIO} & \textbf{Pin Name} & \textbf{Function 0} & \textbf{Function 1} & \textbf{Function 2} & \textbf{Function 3} & \textbf{DRV} & \textbf{Reset} & \textbf{Notes} \\
        \hline
        \endhead
        \endfoot

0 & GPIO0            & GPIO0         & GPIO0   & FSPIQ        & ---            & 2           & 0 & --- \\ \hline
1 & GPIO1            & GPIO1         & GPIO1   & FSPICS0      & ---            & 2           & 0 & --- \\ \hline
2 & MTMS             & MTMS          & GPIO2   & FSPIWP       & ---            & 2           & 1 & --- \\ \hline
3 & MTDO             & MTDO          & GPIO3   & FSPIHD       & ---            & 2           & 1 & --- \\ \hline
4 & MTCK             & MTCK          & GPIO4   & FSPICLK      & ---            & 2           & 1* & --- \\ \hline
5 & MTDI             & MTDI          & GPIO5   & FSPID        & ---            & 2           & 1 & --- \\ \hline
8 & GPIO8            & GPIO8         & GPIO8   & ---          & ---            & 2           & 1 & R \\ \hline
9 & GPIO9            & GPIO9         & GPIO9   & ---          & ---            & 2           & 3 & R \\ \hline
10 & GPIO10          & GPIO10        & GPIO10  & ---          & ---            & 2           & 0 & R \\ \hline
11 & GPIO11          & GPIO11        & GPIO11  & ---          & ---            & 2           & 0 & R \\ \hline
12 & GPIO12          & GPIO12        & GPIO12  & ---          & ---            & 2           & 0 & R \\ \hline
13 & XTAL\_32K\_P    & GPIO13        & GPIO13  & ---          & ---            & 2           & 0 & R \\ \hline
14 & XTAL\_32K\_N    & GPIO14        & GPIO14  & ---          & ---            & 2           & 0 & R \\ \hline
22 & GPIO22          & GPIO22        & GPIO22  & ---          & ---            & 2           & 0 & --- \\ \hline
23 & U0RXD           & U0RXD         & GPIO23  & FSPICS1      & ---            & 2           & 3 & --- \\ \hline
24 & U0TXD           & U0TXD         & GPIO24  & FSPICS2      & ---            & 2           & 4 & --- \\ \hline
25 & GPIO25          & GPIO25        & GPIO25  & FSPICS3      & ---            & 2           & 1 & --- \\ \hline
26 & GPIO26          & GPIO26        & GPIO26  & FSPICS4      & ---            & 3           & 1 & USB \\ \hline
27 & GPIO27          & GPIO27        & GPIO27  & FSPICS5      & ---            & 3           & 3* & USB \\ \hline


    \end{longtable}
\normalsize

\textbf{Drive Strength}

``DRV" column shows the drive strength of each pin after reset:

\begin{itemize}
\item \textbf{0} - Drive current = \verb+~+5 mA
\item \textbf{1} - Drive current = \verb+~+10 mA
\item \textbf{2} - Drive current = \verb+~+20 mA
\item \textbf{3} - Drive current = \verb+~+40 mA
\end{itemize}

\textbf{Reset Configurations}

``Reset" column shows the default configuration of each pin after reset:

\begin{itemize}
\item \textbf{0} - IE = 0 (input disabled)
\item \textbf{1} - IE = 1 (input enabled)
\item \textbf{2} - IE = 1, WPD = 1 (input enabled, pull-down resistor enabled)
\item \textbf{3} - IE = 1, WPU = 1 (input enabled, pull-up resistor enabled)
\item \textbf{4} - OE = 1, WPU = 1 (output enabled, pull-up resistor enabled)
\item \textbf{1*} - If \hyperref[fielddesc:EFUSEDISPADJTAG]{EFUSE\_DIS\_PAD\_JTAG} = 1, the pin MTCK is left floating after reset, i.e., IE = 1. If \hyperref[fielddesc:EFUSEDISPADJTAG]{EFUSE\_DIS\_PAD\_JTAG} = 0, the pin MTCK is connected to internal pull-up resistor, i.e., IE = 1, WPU = 1.
\item \textbf{3*} - IE = 1, WPU = 0. The default value of GPIO27's USB pull-up is 1, which means the pull-up resistor is enabled. For details, please refer to the note below.
\end{itemize}

\textbf{GPIO Input Mode}

The input function of GPIO can be configured as hysteresis or normal mode:
   \begin{itemize}
        \item Hysteresis mode: In the hysteresis mode, the threshold voltage for flipping between high and low levels of GPIO input depends on the direction of level flipping. Specifically, the voltage threshold for flipping from high to low level is slightly lower than the voltage threshold for flipping from low to high level. For details, see Chapter \ref{PadHysEn}. 
        \item Normal mode: Disable hysteresis for GPIO pins, and the threshold voltage for flipping between high and low levels of GPIO input is independent of the direction of level flipping. In other words, the voltage threshold for flipping from high to low level is the same as the voltage threshold for flipping from low to high level.
   \end{itemize}

\textbf{Note:}

\begin{itemize}
\item \textbf{R} - LP pins. Some LP pins have analog functions. For details, see \ref{tab:iomuxgpio-iomax-pins-analog-functions}.
\item \textbf{USB} - USB pull-up resistor enabled 
\begin{itemize}
    \item By default, the USB function is enabled for USB pins (i.e., GPIO26 and GPIO27), and the pin pull-up is decided by the USB pull-up. The USB pull-up is controlled by USB\_SERIAL\_JTAG\_DP/DM\_PULLUP and the pull-up resistor value is controlled by USB\_SERIAL\_JTAG\_PULLUP\_VALUE. For details, see \doclinktrmEN{usbserialjtag} > Chapter \textit{USB Serial/JTAG Controller}.
    \item When the USB function is disabled, USB pins are used as regular GPIOs and the pin's internal weak pull-up and pull-down resistors are disabled by default (configurable by IO\_MUX\_GPIO\regindex{n}\_MCU\_WPU/WPD).
\end{itemize}
\end{itemize}
